\conclusions
 
La caracterización del flujo de subida y renderizado de videos en Picta permitió constatar que el
mecanismo monolítico preexistente constituía el principal cuello de botella del proceso de
publicación. Cualquier interrupción de red durante la transferencia obligaba al administrador de
canal a reiniciar la subida completa desde cero, el servidor Django actuaba como intermediario
innecesario de todos los bytes transferidos, y el procesamiento de metadatos bloqueaba la interfaz
hasta su finalización, impidiendo la gestión concurrente de múltiples contenidos. El estudio de
plataformas con funcionalidades similares confirmó que el protocolo de subida multipart con
persistencia de estado y procesamiento asíncrono constituye el estándar de la industria en sistemas
de gestión de contenido audiovisual, y fundamentó las decisiones de diseño e implementación
adoptadas a lo largo del proyecto.
 
El módulo desarrollado resuelve las deficiencias identificadas mediante tres mecanismos
complementarios. En primer lugar, el sistema de carga fragmentada divide el archivo en fragmentos
de tamaño configurable y persiste el índice del último fragmento confirmado tanto en IndexedDB del
navegador como en el servidor, garantizando que ante cualquier forma de interrupción ---voluntaria,
accidental o por fallo de red--- la transferencia pueda retomarse desde el punto exacto de
interrupción sin retransmitir datos ya enviados. En segundo lugar, la delegación de la transferencia
de bytes directamente al almacenamiento MinIO S3 mediante URLs prefirmadas eliminó al servidor
Django como proxy de la transferencia, reduciendo su carga de CPU y consumo de ancho de banda
durante las subidas. En tercer lugar, la introducción del estado intermedio \textit{sin\_procesar}
y la delegación del procesamiento de metadatos a una cola de tareas Celery desacoplaron
completamente la subida del renderizado final, permitiendo al administrador de canal gestionar
múltiples contenidos de forma concurrente sin esperas bloqueantes.
 
La evaluación del sistema, realizada mediante trece casos de prueba funcionales y diez pruebas de
aceptación contrastadas contra los criterios de las historias de usuario, arrojó un índice de
cumplimiento del 100~\% en cuatro navegadores de escritorio. Las métricas de tiempo de subida
resultaron proporcionales al tamaño del archivo sin degradación apreciable al aumentar el número de
fragmentos, y la comparación con el sistema anterior confirmó mejoras concretas en tolerancia a
fallos, persistencia de progreso, control del usuario, rendimiento del servidor y gestión
concurrente de sesiones. Estos resultados validan que las decisiones arquitectónicas adoptadas
se traducen en un avance técnico sustancial respecto al mecanismo que operaba previamente en
producción.